\documentclass[conference]{IEEEtran}
\IEEEoverridecommandlockouts

% 中文会议稿：需 XeLaTeX / LuaLaTeX + ctex
\usepackage{ctex}
\usepackage{cite}
\usepackage{amsmath,amssymb,amsfonts}
\usepackage{algorithm}
\usepackage{algpseudocode}
\usepackage{graphicx}
\usepackage{booktabs}
\usepackage{xcolor}
\usepackage{tikz}
\usetikzlibrary{arrows.meta,positioning,fit,shapes.geometric}
\usepackage[hidelinks]{hyperref}
\usepackage{url}
\usepackage{multirow}

\newcommand{\code}[1]{\texttt{\small #1}}

\definecolor{boxblue}{RGB}{225,235,250}
\definecolor{boxgreen}{RGB}{224,242,228}
\definecolor{boxorange}{RGB}{252,238,220}
\definecolor{boxgray}{RGB}{238,238,238}

\begin{document}

\title{面向 LLM 安全情报采集的混合控制智能体：\\可审计预算约束下的自主检索与无强化学习决策增强}

\author{\IEEEauthorblockN{作者待补}
\IEEEauthorblockA{\textit{单位待补}\\
\textit{邮箱待补}}
}

\maketitle

\begin{abstract}
大语言模型（LLM）应用安全威胁证据分散于 NVD、CISA KEV、OSV 与 arXiv 等异构公开源。开放式采集若仅依赖固定源轮询与模板检索，难以在有限 API 预算下兼顾覆盖完整性与噪声控制。本文报告一套面向 LLM 安全领域的情报采集智能体系统：在 Claude Agent SDK 原生工具循环中赋予智能体检索词生成与高级算子组合能力，同时由确定性 Python 控制面保留轮次骨架、预算安全阀、UCB 组合审批与持久化，形成\textbf{混合控制面}。在不引入 Agentic 强化学习的前提下，系统以 UCB 风格效用对智能体提交的查询候选做审批门控，并以自然语言技巧库（Playbook）跨运行沉淀高收益算子组合。离线回放 33 次历史运行、617 轮查询结果表明，相对全源轮询基线，自适应 UCB 的归一化 regret 由 1.245 降至 1.063；端点能力矩阵上工具调用与结构化决策解析成功率均达 100\%。本文贡献在于给出可审计、可降级的自主采集架构，并明确无 RL 条件下检索决策增强的工程边界。
\end{abstract}

\begin{IEEEkeywords}
大语言模型安全, 网络威胁情报, 自主智能体, 混合控制, 多臂老虎机, 工具调用, 可审计性
\end{IEEEkeywords}

% ---------- English abstract (bilingual) ----------
\section*{Abstract (English)}
Evidence of LLM-application security threats is scattered across heterogeneous public feeds such as NVD, CISA KEV, OSV, and arXiv. Open-ended collection with fixed source rotation and template queries struggles to balance coverage and noise under a finite API budget. We present an intelligence-collection agent for the LLM-security domain that grants free-form query generation and advanced operator composition inside a Claude Agent SDK tool loop, while a deterministic Python control plane retains the round skeleton, budget safety valves, UCB portfolio approval, and persistence---a \emph{hybrid control plane}. Without Agentic reinforcement learning, the system gates agent-proposed query candidates with a UCB-style utility and accumulates high-yield operator recipes in a natural-language Playbook across runs. Offline replay over 33 historical runs and 617 rounds shows that adaptive UCB reduces normalized regret from 1.245 to 1.063 versus round-robin; an endpoint capability matrix reports 100\% success for both tool calling and structured decision parsing. We contribute an auditable, degradable autonomous collection architecture and clarify the engineering boundary of retrieval-decision improvement without RL.

\section{引言}
\label{sec:intro}

LLM 正以对话助手、检索增强生成（RAG）应用与工具调用智能体等形式进入生产系统。随之而来的是一类相对独立的安全问题：提示注入、越狱、RAG 投毒、智能体工具滥用、模型供应链风险等，已被 OWASP LLM Top~10 系统化~\cite{owasp2025llmtop10}，并在真实集成场景中得到实证~\cite{greshake2023injection}。与单一 CVE 订阅不同，这类威胁的证据高度碎片化：可利用漏洞常先出现在 NVD~\cite{nvd} 与 CISA KEV~\cite{cisakev}；受影响开源组件出现在 OSV~\cite{osv}；攻击手法与防御评估则往往更早出现在 arXiv 预印本~\cite{arxivapi}。因此，面向 LLM 安全的情报采集天然是\textbf{多源、多模态、预算受限}的序贯决策问题。

现有流水线常见做法是：固定源集合轮询、按话题关键词表生成模板查询、由用户指定最大轮数。黑板（blackboard）上虽已记录 novelty、duplicate ratio、noise ratio 与覆盖缺口，这些信号却很少真正进入决策者。与此同时，完全放开智能体自主调用外部 API 又会带来预算失控、重复查询与不可审计风险。Agentic AI 在网络安全中的机遇与治理缺口已被广泛讨论~\cite{kshetri2025agentic,lazer2026survey}，但针对\textbf{LLM 安全情报采集}这一具体任务，仍缺少可落地的、把自主性关进可审计控制面的系统报告。

本文聚焦两个相互耦合的研究问题：
\begin{itemize}
  \item \textbf{RQ1（混合控制）}：如何在赋予智能体轮内检索自主性的同时，用确定性控制面保证预算、去重与终止条件可审计、可强制执行？
  \item \textbf{RQ2（无 RL 决策增强）}：在不训练策略网络的前提下，如何用 UCB 审批门控与跨运行 Playbook，提升源/查询/算子选择的长期收益？
\end{itemize}

我们的回答是一套已落地的自包含引擎（\texttt{backend/intel\_agent}）：智能体在 MCP 白名单工具上自由组合高级检索算子~\cite{anthropicmcp,claudeagentsdk}；控制面以 hooks、UCB 组合选择与多层 fallback 约束其行为~\cite{claudeagenthooks}；Playbook 以自然语言技巧形式沉淀可复用经验。下游知识图谱构建（如 CTINexus~\cite{cheng2025ctinexus}）与采集解耦，失败互不阻断。

\noindent\textbf{贡献概括：}
\begin{enumerate}
  \item 提出并实现面向 LLM 安全情报采集的\textbf{混合控制面}架构：SDK 工具循环负责探索，Python 控制器负责预算、审批、覆盖与落盘。
  \item 给出\textbf{无 RL} 的检索决策增强路径：智能体提案 $\rightarrow$ UCB 效用审批 $\rightarrow$ 归一化奖励回写 $\rightarrow$ Playbook 召回/衰减。
  \item 报告可复核的子系统证据：UCB 离线 regret 对比、端点能力矩阵，以及与确定性规则基线共享的降级契约。
\end{enumerate}

\section{相关工作}
\label{sec:related}

\subsection{LLM 安全与威胁情报}
提示注入与间接注入揭示了“数据即指令”的新攻击面~\cite{greshake2023injection}；OWASP 持续维护 LLM 应用风险清单~\cite{owasp2025llmtop10}。传统 CTI 强调 STIX/TAXII 结构化交换与 ATT\&CK 知识组织，而 LLM 安全证据大量仍以漏洞公告与学术摘要形式存在。CTINexus 展示了用 LLM 上下文学习从非结构化 CTI 构建知识图谱的路径~\cite{cheng2025ctinexus}，但其输入假设情报文本已可得。本文关注其上游：如何在预算约束下从活源\textbf{采到}可用情报。

\subsection{工具调用智能体}
ReAct 将推理轨迹与环境动作交错，使 LLM 能基于观察修正计划~\cite{yao2023react}；Toolformer 表明模型可自学何时调用工具~\cite{schick2023toolformer}。Claude Agent SDK 进一步提供进程内自定义工具、hooks 与会话管理~\cite{claudeagentsdk,claudeagenthooks}。本文不把智能体当作端到端黑盒执行器，而是将其嵌入混合控制面：轮内动作可自主，轮间状态与安全阀由确定性代码掌握。

\subsection{检索中的多臂老虎机}
信息检索中，多查询可被同时保持为活跃臂，由 bandit 在速率限制下分配拉取轮次~\cite{li2016multiple}；文档判定与主题查询识别也可建模为探索—利用问题~\cite{losada2016feeling,losada2021cost}。UCB1 提供有限时间 regret 保证且实现简单~\cite{auer2002ucb}。与直接让 bandit 执行查询不同，本文将 UCB 用作\textbf{审批门}：臂由智能体提出（含高级算子签名），确定性评分决定本轮至多批准的调用组合，从而保留语义生成能力又限制工具滥用。

\subsection{Agentic 网络安全}
近期综述与政策讨论强调自主智能体在 SOC 自动化中的潜力及其治理缺口~\cite{kshetri2025agentic,lazer2026survey}。本文贡献是领域收窄后的可运行架构：固定威胁话题分类（Core/Extended）、四源 typed 工具、可观测黑板与离线可回放的决策日志。

\section{问题设定与设计目标}
\label{sec:problem}

设采集运行由轮次 $t=1,\ldots,T$ 组成，每轮可发起有限次源 API 调用。状态包含：已采集去重条目集合、各 Core 话题覆盖缺口、历史查询收益、剩余 API 预算。动作空间为四元组 $(\textit{source}, \textit{query}, \textit{operators}, \textit{stop})$。设计目标：
\begin{enumerate}
  \item \textbf{自主性}：查询文本与算子组合由智能体在工具循环中生成，而非仅模板改写。
  \item \textbf{可审计预算}：任何调用可追溯至批准的候选 ID；超预算、未批准调用被 hooks/处理器拒绝。
  \item \textbf{可降级}：SDK 不可用或结构化决策失败时，行为不低于共享规则基线。
  \item \textbf{可演化}：高收益算子组合跨运行沉淀，低收益技巧衰减。
\end{enumerate}

非目标：本文不训练 Agentic RL 策略；不把 item 级知识图谱融合纳入采集主循环。

\section{系统架构：混合控制面}
\label{sec:arch}

\subsection{总体结构}
系统分为四层（图~\ref{fig:arch}）：（1）\textbf{模式层}定义全量（full）与增量（incremental）目标；（2）\textbf{控制面}（\texttt{IntelAgentController}）编排轮次、相关性标注、缺口重算、终止与持久化；（3）\textbf{智能体面}（\texttt{IntelAgentLoop}）在 SDK 会话中调用 MCP 工具；（4）\textbf{确定性增强}含 UCB、Playbook、三层相关性与规则 fallback。

\begin{figure}[!t]
\centering
\begin{tikzpicture}[
  font=\scriptsize,
  box/.style={draw, rounded corners=2pt, align=center, inner sep=4pt},
  arr/.style={-{Latex}, thick}
]
\node[box, fill=boxblue, minimum width=7.2cm] (mode) {CollectionMode\\full 配额 / incremental CorpusGap};
\node[box, fill=boxorange, minimum width=7.2cm, below=0.35cm of mode] (ctrl) {控制面 Controller\\预算 $\cdot$ UCB 审批 $\cdot$ 终止 $\cdot$ 落盘};
\node[box, fill=boxgreen, minimum width=3.3cm, below left=0.45cm and -0.1cm of ctrl] (agent) {SDK Agent Loop\\typed 源工具};
\node[box, fill=boxgray, minimum width=3.3cm, below right=0.45cm and -0.1cm of ctrl] (det) {UCB / Playbook\\Relevance / Rules};
\node[box, fill=white, minimum width=7.2cm, below=1.55cm of ctrl] (store) {Blackboard + Trace JSONL + Mongo/JSON Store};
\draw[arr] (mode) -- (ctrl);
\draw[arr] (ctrl) -- (agent);
\draw[arr] (ctrl) -- (det);
\draw[arr] (agent) -- (store);
\draw[arr] (det) -- (store);
\draw[arr] (ctrl) -- (store);
\end{tikzpicture}
\caption{混合控制面架构：智能体负责轮内检索探索，控制面与确定性模块负责审批、度量与持久化。}
\label{fig:arch}
\end{figure}

\subsection{单轮协议}
每轮执行算法~\ref{alg:round}。关键不变量是：\textbf{智能体可以提案，但不能绕过审批与预算}。自适应增量模式下，源工具调用必须携带已批准的 \code{candidate\_id}；PreToolUse hook 在预算耗尽或未批准时拒绝调用~\cite{claudeagenthooks}。

\begin{algorithm}[!t]
\caption{混合控制面单轮协议}
\label{alg:round}
\begin{algorithmic}[1]
\State 召回 Playbook 技巧；渲染黑板摘要 digest（禁止倾倒全文）
\State 启动/恢复 SDK 会话；智能体调用源工具并提交本轮摘要
\State \textbf{若} 无有效采集 \textbf{则} 运行规则 fallback 轮次
\State 三层相关性标注；合并入 blackboard；harvest 高收益技巧
\State 重算 CorpusGap / RunGap；记录每调用归一化奖励
\State Critic + 饱和/停滞/预算规则决定是否继续
\end{algorithmic}
\end{algorithm}

\subsection{安全边界与工具白名单}
进程内 MCP 仅暴露源搜索、状态查询、候选评分、Playbook 与轮次摘要工具；禁用 Bash/文件系统类内建工具，避免“采集智能体”演变为通用代码代理。Docker 以 root 运行时需显式沙箱标志，否则权限模式会静默降级。该设计把\textbf{能力面}与\textbf{权限面}分离：自主性体现在检索策略，而非主机控制。

\subsection{多层 Fallback}
\begin{itemize}
  \item \textbf{决策层}：结构化终止/计划解析失败 $\rightarrow$ 规则裁决；
  \item \textbf{轮次层}：SDK 会话失败且无条目 $\rightarrow$ 确定性扫源；
  \item \textbf{策略层}：自适应初始化失败 $\rightarrow$ legacy 配额路径；
  \item \textbf{相关性层}：embedding/LLM 失败 $\rightarrow$ 规则分。
\end{itemize}
所有降级写入 blackboard.errors 与 trace，避免静默失败。

\section{无 RL 的检索决策增强}
\label{sec:decision}

\subsection{智能体提案与 UCB 审批}
自适应增量中，智能体通过 \code{score\_query\_candidates} 提交 $1$--$8$ 个候选。臂键为
\[
\textit{topic}\mid\textit{channel}\mid\textit{source}\mid\textit{intent}\mid\textit{operator\_signature}.
\]
效用近似为
\[
u = \widehat{r} - 0.2\,\textit{noise} - 0.2\,\textit{dup} + \sqrt{\frac{2\ln(n+2)}{n_a+1}},
\]
其中 $\widehat{r}$ 融合目标缺口优先级与历史均值奖励，末项为 UCB 探索奖励~\cite{auer2002ucb}。\code{select\_portfolio} 每轮最多批准 4 次调用，并对 discovery/confirmation 意图施加预算占比上限。该机制把 bandit 从“替智能体选动作”改为“对智能体提案做资源分配”，更贴近 IR 中在速率限制下调度多个查询的设定~\cite{li2016multiple,losada2021cost}。

\subsection{归一化奖励}
单次调用奖励由缺口优先级下降、discovery/confirmation 加成，以及噪声/重复惩罚组成，裁剪到 $[-1,1]$。奖励进入 \code{query\_outcomes}，既服务在线 UCB，也支持离线回放对比。

\subsection{Playbook 自演化}
高产出调用（例如每调用新增相关条目 $\ge 0.5$）被 harvest 为自然语言技巧，键为 $\textit{topic\_bucket}::\textit{source}::\textit{operator\_signature}$。下轮按语义相似与历史收益召回 Top-$k$ 注入 digest；长期低收益条目衰减并可能废弃。与参数化策略梯度不同，Playbook 是\textbf{可人工审阅}的经验记忆，适合安全运营场景对可解释性的要求。

\subsection{高级检索算子}
typed 工具直接暴露各源算子：NVD 的 CWE/CVSS/时间窗/KEV；arXiv 的布尔与分类限定；OSV 的生态/包/purl；CISA KEV 的 CVE 列表。算子签名进入臂键，使“会用高级语法”的查询与裸关键词查询在统计上可区分。

\subsection{为何暂不做 Agentic RL}
当前主流托管模型不可本地训练；相关性金标尚未规模化，奖励噪声会主导策略更新；而上下文工程、UCB 门控与 Playbook 的边际收益仍未耗尽。本文将 RL 留作评估基线成熟后的重估选项。

\section{相关性、覆盖与终止}
\label{sec:relevance}

三层过滤（规则 $\rightarrow$ embedding $\rightarrow$ flash LLM）把条目标为相关/不相关/不确定，并写入 \code{method} 字段以便分层评估。Full 模式以 Core 话题配额驱动覆盖；Incremental adaptive 以历史语料上的 CorpusGap（话题 $\times$ 证据渠道）驱动软覆盖。终止由 critic 解释，但硬约束顺序为：\code{max\_rounds} / API 预算 / 超时 $\rightarrow$ 饱和或停滞 $\rightarrow$ \code{min\_rounds} 下限。自适应路径在连续低收益且最大 UCB 效用过低时标记缺口耗尽并停止，避免在不可达主题上空转。

\section{实现与运行形态}
\label{sec:impl}

系统以 Python 包自包含交付，CLI 提供 \code{full}/\code{incremental}/\code{latest}。可选 MongoDB 以 \code{item\_id} 全局去重并服务 CorpusGap；\code{--verbose} 将工具调用与 critic 决策落入 JSONL，供控制台实时订阅。采集与 \texttt{ctinexus\_kg} 禁止互相 import，仅通过 \code{run\_id} 与持久化条目交接。

\section{评估证据}
\label{sec:eval}

\subsection{端点能力矩阵}
在 DeepSeek Anthropic 兼容端点上，工具调用成功率 $20/20$，tool-forcing 结构化决策解析 $20/20$（目标分别为 $\ge 95\%$ 与 $\ge 90\%$）。子智能体与提示缓存未达标，故终止 critic 采用独立快速结构化调用，上下文靠摘要压缩而非缓存。

\subsection{UCB 离线回放}
对历史 blackboard 中的 \code{query\_outcomes} 回放（33 次运行、617 轮），自适应 UCB 归一化 regret 为 \textbf{1.063}，全源轮询为 \textbf{1.245}，满足“不低于轮询”的验收门槛。该实验不消耗在线 API，可重复。

\subsection{在线冒烟}
一次全量冒烟在 API 预算用尽前采集 255 条，多数 Core 话题达到配额，残留缺口可观测。这证明控制面预算阀与覆盖记账可用，但\textbf{不能}替代同预算 A/B 与人工精确率抽检。

\subsection{局限}
完整同预算 A/B（冻结目标 $\times$ 引擎 $\times$ 重复）与人工金标精确率仍在进行中；因此本文不宣称端到端“智能提升百分比”，而将 UCB regret 与端点可靠性作为子系统级证据。相关性规则层仍可能偏乐观，需依赖第 2--3 层与人工校准。

\section{讨论}
\label{sec:discuss}

\textbf{RQ1。}混合控制面把“自主”限定在检索策略空间：智能体可生成查询与算子，但不能突破预算、白名单与审批。hooks + 候选 ID + 共享规则 fallback 使每次降级可审计，这与纯 ReAct 式端到端代理形成对照~\cite{yao2023react}。

\textbf{RQ2。}UCB 审批保留探索项，同时用缺口优先级引导利用；Playbook 提供跨运行记忆而无须可微策略。该组合回答了“无 RL 还能提升什么”：提升的是\textbf{资源分配与经验复用}，不是端到端策略最优性。

对运营的启示是：开放采集应先定义覆盖与预算目标，再谈智能体；否则无法区分“更聪明”与“更会烧配额”。

\section{结论与展望}
\label{sec:concl}

本文报告了面向 LLM 安全情报采集的混合控制智能体：以 SDK 工具循环实现轮内自主检索，以确定性控制面保证可审计预算与可降级行为，并以 UCB 审批与 Playbook 在无 RL 条件下增强决策。后续工作包括完成同预算 A/B 与人工精确率协议、强化增量水位线在 Mongo 模式下的端到端验证，以及将高置信条目接入知识图谱融合流水线。

\section*{致谢}
作者信息与资助声明待补。

\section*{数据与代码可用性}
系统实现位于项目仓库 \texttt{backend/intel\_agent/}；UCB 回放脚本为 \texttt{scripts/bandit\_replay.py}；运行轨迹可导出为 JSONL。活源 API 密钥与原始情报文本因许可与敏感原因不公开完整转储。

\section*{作者贡献（CRediT）}
待作者名单确定后按 CRediT 十四项角色填写。

\section*{利益冲突}
作者声明无利益冲突。

\section*{AI 使用披露}
本文初稿由学术写作辅助智能体基于系统代码与内部技术文档生成，经作者审阅、核实引用与实验数字后定稿。所有参考文献均要求可通过 DOI 或官方文档核验；未核验条目不得保留。

\bibliographystyle{IEEEtran}
\bibliography{references}

\end{document}
